\documentclass[12pt]{article}
\usepackage[margin=0.58in,top=0.62in,bottom=0.48in]{geometry}
\usepackage{lmodern}
\usepackage{microtype}
\usepackage{fancyhdr}
\usepackage{titlesec}
\usepackage{enumitem}
\usepackage{lastpage}
\usepackage[hidelinks]{hyperref}
\setlength{\parindent}{0pt}
\setlength{\parskip}{0.18em}
\setlength{\headheight}{26pt}
\raggedbottom
\titleformat{\section}{\normalsize\bfseries\scshape}{}{0pt}{}
\titleformat{\subsection}{\normalsize\bfseries}{}{0pt}{}
\titlespacing*{\section}{0pt}{0.35em}{0.12em}
\titlespacing*{\subsection}{0pt}{0.25em}{0.08em}
\pagestyle{fancy}
\fancyhf{}
\fancyhead[C]{\footnotesize Lit Example Reader \quad Assignment ID: U1L18LE \quad Page \thepage\ of \pageref{LastPage}}
\renewcommand{\headrulewidth}{0.5pt}
\newenvironment{playpassage}{\begin{quote}\footnotesize\setlength{\parskip}{0.08em}\setlength{\parindent}{0pt}}{\end{quote}}
\newcommand{\speaker}[1]{\textbf{#1.}}
\begin{document}
\begin{center}
{\LARGE\bfseries Week 18 Lit Example Reader}\par\vspace{0.02em}
{\large\scshape Revising a Shakespeare Judgment}\par\vspace{0.06em}
\rule{0.78\textwidth}{0.5pt}
\end{center}

\textbf{Purpose:} Practice a visible revision trail on familiar Shakespeare evidence before revising a portfolio artifact of your own.

\textbf{What to watch for:} A revision should preserve the original claim, diagnose one exact weakness, change the reasoning, and explain why the change matters.

\section*{Before the Passages: The Week 18 Revision Tool}
\begin{itemize}[leftmargin=1.3em,itemsep=0.05em,topsep=0.08em,parsep=0.03em]
\item \textbf{Recover:} identify the draft's claim, evidence, inference, and assumption.
\item \textbf{Diagnose:} apply one fitting criterion to a precise feature.
\item \textbf{Revise:} clarify, narrow, add support, or repair the connection.
\item \textbf{Caption:} name the before-and-after change and the habit it shows.
\end{itemize}

\section*{Passage 1: Macbeth, Act 4, Scene 1 --- Security from a Promise}
\textbf{Context:} Macbeth hears the second apparition promise that none ``of woman born'' shall harm him. He turns the promise into confidence about Macduff.
\begin{playpassage}
\speaker{Second Apparition} Be bloody, bold, and resolute; laugh to scorn\\
The power of man, for none of woman born\\
Shall harm Macbeth.\\[0.25em]
\speaker{Macbeth} Then live, Macduff: what need I fear of thee?\\
But yet I'll make assurance double sure,\\
And take a bond of fate: thou shalt not live.
\end{playpassage}

\subsection*{Assisted Revision Model}
Suppose an early caption said: ``Macbeth is fooled because the witches lie.'' That claim notices danger but outruns the passage. The words are equivocal, yet the later plot gives them a technical fulfillment. A stronger diagnosis is that the caption treats deception as a simple false statement and misses Macbeth's inference from an ambiguous promise to practical certainty. A substantive revision would narrow the claim: Macbeth accepts the most reassuring sense, fails to test the source and terms, and then uses that confidence to justify further violence. The revision is stronger because it names evidence, a logical habit, and the exact point of failure.

\textbf{Reading task:} Underline the promise; circle Macbeth's inference; bracket the action he adds. In the model, mark the original claim, diagnosis, and revised claim.

\vfill
\section*{Annotation Notes}
\noindent\rule{\textwidth}{0.4pt}\\[1.0em]
\rule{\textwidth}{0.4pt}\\[1.0em]
\rule{\textwidth}{0.4pt}

\newpage
\section*{Passage 2: Julius Caesar, Act 2, Scene 1 --- Possibility as Necessity}
\textbf{Context:} Brutus has no present personal cause against Caesar. He imagines how a crown might change Caesar and reaches a preventive conclusion.
\begin{playpassage}
\speaker{Brutus} It must be by his death: and for my part,\\
I know no personal cause to spurn at him,\\
But for the general. He would be crown'd:\\
How that might change his nature, there's the question.\\
...\\
And therefore think him as a serpent's egg\\
Which, hatch'd, would, as his kind, grow mischievous,\\
And kill him in the shell.
\end{playpassage}

\subsection*{Reader Notes}
This passage shifts the revision work to you. An early response might praise Brutus's concern for Rome or condemn him as simply wrong. Neither judgment is yet a logic analysis. Attend to the distance between \textit{might change} and \textit{must be by his death}, the work performed by the serpent's-egg analogy, and the alternatives or evidence left unexamined. Decide which criterion best diagnoses the early response; do not merely attach every course term.

\textbf{Reading task:} Underline the uncertain language. Box the conclusion of necessity. Star the analogy. In the margin, name one kind of support the argument still needs.

\section*{How to Use These Passages on the Worksheet}
Passage 1 demonstrates the shape of a revision trail. For Passage 2, preserve the early claim, diagnose one exact weakness, revise it with textual evidence, and write a caption that explains what the later Whately tool made visible.

\vfill
\section*{Transfer Notes}
\noindent\rule{\textwidth}{0.4pt}\\[1.0em]
\rule{\textwidth}{0.4pt}\\[1.0em]
\rule{\textwidth}{0.4pt}\\[1.0em]
\rule{\textwidth}{0.4pt}\\[1.0em]
\rule{\textwidth}{0.4pt}
\end{document}
